\section{All sectors and complete rationality}\label{sec:exhaustion}

The tensor functor constructed in Theorem~\ref{thm:cwx-embedding} gives
finite-index representations of the net.  We must still show that it accounts
for every finite sector.  We prove the stronger statement that every
irreducible locally normal representation comes from an ordinary unitary
module.  The point is to reconstruct a weak module before assuming that the
representation has finite energy multiplicities.  Regularity will then supply
those multiplicities.

The scalar VOA has already been treated separately.  We first consider
\(V\ne\mathbb C\Omega\).  All representations of \(A_V\) are understood
to be compatible families of normal representations of its interval algebras.

\begin{proposition}\label{prop:exhaustion}
Let \(V\ne\mathbb C\Omega\) be a simple unitary strongly rational VOA
which is energy-bounded and strongly local.  Every nonzero irreducible
locally normal representation
\((\mathcal K,\pi)\) of \(A_V\) is the strong integration of an ordinary
simple unitary \(V\)-module.  No finite-index or finite-energy-multiplicity
assumption on \(\pi\) is required.
\end{proposition}

\begin{proof}
Write \(A=A_V\).  Reconstruction will give a weak module on the finite-energy vectors.
We first establish the separability and discrete rotation spectrum
needed for that construction; regularity will then produce a simple
module whose Hilbert closure is the whole sector.  The vacuum Hilbert space is separable, since it is
the completion of the countable sum of the finite-dimensional spaces
\(V_n\).  The Hilbert space \(\mathcal K\) is also separable.  To see this,
choose a countable basis of intervals and, in the unit ball of each
corresponding vacuum algebra, a countable \(\sigma\)-strong-* dense subset.
Such subsets exist because these von Neumann algebras have separable
preduals.  Normality and additivity imply that the images of these countably
many subsets generate all the represented local algebras.  The norm-closed
unital *-algebra they generate is separable and acts irreducibly.  Every
nonzero vector is therefore cyclic for this algebra, which proves the claim.

The canonical inner M\"obius implementation in a locally normal
representation has positive energy
\cite[Theorem~3.8]{Weiner}; its Virasoro action and the associated annulus
action are the ones used in \cite[Sections~2, 4 and 9]{HTVOA}.
Let \(H\geq0\) be the self-adjoint generator of the lifted rotation group.
Covariance implies that \(e^{2\pi iH}\) commutes with every represented
local algebra.  Irreducibility makes it scalar.  Thus, for some
\(h\in[0,1)\), the spectral theorem gives
\[
  \operatorname{spec}H\subset (h+\mathbb Z)\cap[0,\infty).
\]
In particular, \(H\) has pure point spectrum.  Its eigenspaces
\(\mathcal K_\lambda=\ker(H-\lambda)\) need not yet be finite dimensional.
Set
\[
  M=\bigoplus_{\lambda\in\operatorname{spec}H}^{\mathrm{alg}}
        \mathcal K_\lambda.
\]
This is a dense subspace of \(\mathcal K\).

The VOA reconstructed from \(A_V\) is \(V\), with its given conformal
vector and inner product \cite[Theorem~8.9(ii)]{HTVOA}.  Here strong
locality supplies AQFT-locality in the terminology of that theorem, and
the required finite-dimensional vacuum energy spaces are precisely the
spaces \(V_n\).  We now apply Proposition~9.3, Definition~9.4, Theorem~9.5 and
Corollary~9.6 of \cite{HTVOA}, together with Remark~9.8 of that paper.
The remark extends all these conclusions to discrete rotation spectra
with possibly infinite-dimensional eigenspaces.  They provide a weak
unitary \(V\)-module structure \(Y^M\) on \(M\) such that
\begin{enumerate}[label=\textup{(\roman*)}]
  \item the conformal operator \(L_0^M\) equals \(H|_M\);
  \item for every homogeneous \(a\in V\), smooth test function \(f\),
        and interval \(I\) containing its support,
        \begin{equation}\label{eq:sector-closed-fields}
          \pi_I\!\left(\overline{Y(a,f)}\right)
             =\overline{Y^M(a,f)|_M}.
        \end{equation}
\end{enumerate}
On the left, a normal representation acts on an affiliated closed operator
through its spectral and polar data.  In particular, \eqref{eq:sector-closed-fields}
asserts equality of closed operators, and \(M\) is a core for the operator
on either side.  It does not assert only equality of formal matrix
coefficients.  This is the part of the reconstruction theorem needed below.

Strong rationality implies regularity: every weak \(V\)-module is an
algebraic direct sum of ordinary simple modules
\cite[Theorem~4.5]{ABD}.  Choose a
nonzero ordinary simple submodule \(S\subset M\).  Its inherited
Hermitian form is positive and invariant.  Because the two conformal
operators agree, the homogeneous spaces
\[
   S_\lambda=S\cap\mathcal K_\lambda
\]
are finite dimensional.  Let \(Q\) be the orthogonal projection of
\(\mathcal K\) onto the Hilbert closure of \(S\).  Orthogonality of the
energy spaces gives
\[
   \overline S=\mathop{\widehat\bigoplus}_{\lambda}S_\lambda,
   \qquad Q(M)\subset S\subset M.
\]
We will show that \(Q\) commutes with the represented net.  The domain
argument is necessary here: algebraic invariance of \(S\) alone would not
give reduction of the closed smeared fields.

Both each mode and its algebraic adjoint preserve \(S\).  Indeed, the
unitary adjoint formula expresses the latter as a finite linear combination
of modes of PCT-transformed states and their \(L_1\)-descendants.
Consequently \(S^\perp\cap M\) is invariant under every mode, and
\[
  Q a_{(n)}^M x=a_{(n)}^M Qx
       \qquad(a\in V,\ n\in\mathbb Z,\ x\in M).
\]
The projection \(Q\) also commutes with every energy projection.  For
\(x\in M\), the component of \(Y^M(a,f)x\) at a fixed energy is a
finite sum of mode terms, since \(x\) has only finitely many energy
components.  Comparing these components proves
\[
   QY^M(a,f)x=Y^M(a,f)Qx\qquad(x\in M).
\]
Both sides are Hilbert-space vectors by the reconstruction theorem, so
componentwise equality suffices.  The argument extends by finite sums to
inhomogeneous \(a\).

Put \(T=\overline{Y^M(a,f)|_M}\).  If \(x\in\operatorname{Dom}T\),
choose \(x_j\in M\) with \(x_j\to x\) and \(Tx_j\to Tx\).
Then
\[
   Qx_j\longrightarrow Qx,
   \qquad TQx_j=QTx_j\longrightarrow QTx.
\]
Closedness implies \(Qx\in\operatorname{Dom}T\) and \(TQx=QTx\).
The same holds with \(1-Q\) in place of \(Q\).  Hence the orthogonal
decomposition defined by \(Q\) reduces \(T\), and \(Q\) commutes with
its bounded polar data.  By \eqref{eq:sector-closed-fields}, these are the
data of the represented vacuum smear.  Such data generate
\(\pi_I(A(I))\), because vacuum smears generate \(A(I)\) and
\(\pi_I\) is normal.  Thus \(Q\) commutes with every represented local
algebra.

Irreducibility and \(S\ne0\) now force \(Q=1\).  Each \(S_\lambda\)
is closed in \(\mathcal K_\lambda\), so this identity implies
\(S_\lambda=\mathcal K_\lambda\) for every \(\lambda\).  Therefore
\(M=S\).  In particular the energy multiplicities are finite and
\(M\) is an ordinary simple unitary module.  Finally,
\eqref{eq:sector-closed-fields} identifies its strongly integrated
representation with \(\pi\).  The identification uses the actual closed
smeared fields, as required by the CWX construction.
\end{proof}

\begin{theorem}\label{thm:rationality}
Let \(V\) be a simple unitary strongly rational VOA.  Its conformal net
\(A_V\) is completely rational.  The CWX functor constructed in
Theorem~\ref{thm:cwx-embedding} is essentially surjective onto the
additive category of finite-index locally normal DHR representations of
\(A_V\).  Consequently it is a braided unitary tensor equivalence
\[
  F^V_{\mathrm{CWX}}:
  \operatorname{Rep}^{u}(V)\xrightarrow{\ \simeq\ }
  \operatorname{Rep}^{f}(A_V).
\]
\end{theorem}

\begin{proof}
For the nontrivial VOA, Corollary~\ref{cor:locality} supplies the net and
the hypotheses of Proposition~\ref{prop:exhaustion}.  That proposition
shows that every irreducible locally normal sector belongs to the CWX
image.  There are only finitely many such sectors, since \(V\) has only
finitely many simple modules.  They have finite statistical dimension:
Theorem~\ref{thm:cwx-embedding} sends their rigid module duals and duality
maps to net duals and duality maps.  No finite-index assumption was used
in Proposition~\ref{prop:exhaustion} to obtain this conclusion.

The irreducible diffeomorphism-covariant circle net \(A_V\) has the split
property by \cite[Theorem~5.4]{MTW}.  The dichotomy
\cite[Theorem~4.9]{LongoXu} applies to a split conformal net all of whose
irreducible sectors have finite statistical dimension.  It says that such
a net is completely rational or has uncountably many inequivalent
irreducible sectors.  The finite list above excludes the second
alternative.  Thus \(A_V\) is strongly additive and has finite
\(\mu\)-index, as well as being irreducible and split.

For a locally normal representation on a separable Hilbert space,
complete rationality gives a direct sum of irreducible sectors of finite
statistical dimension \cite[Corollary~39 and the proof of Corollary~40; Appendix~B]{KLM}. The modular-PCT hypothesis
is automatic for conformal nets, as explained in Section~3, before
Corollary~10, of \cite{KLM}.

The same conclusion holds without separability. Indeed, apply the
countable-generator construction from Proposition~\ref{prop:exhaustion}
to the present representation $\pi$. The resulting separable unital
C*-algebra $D$ satisfies
\[
 D''=\bigvee_I\pi_I(A_V(I)).
\]
For any vector $\xi$, the separable cyclic space $\overline{D\xi}$
reduces $D$, hence also $D''$ and every represented interval algebra.
A maximal orthogonal family of these cyclic reducing spaces exhausts
the representation Hilbert space: any nonzero orthogonal complement
would supply another such space. Compression to a reducing space
preserves local normality. Apply the separable decomposition
\cite[Corollary~39 and the proof of Corollary~40; Appendix~B]{KLM}
to each restriction and combine the resulting orthogonal sums.

In a dualizable representation only finitely many irreducible summands
can occur. Any finite sub-sum of $n$ nonzero summands has statistical
dimension at least $n$ and at most the dimension of the representation.
This excludes infinite multiplicity even if the full decomposition has
an uncountable index set.  The CWX functor preserves finite direct sums,
and Proposition~\ref{prop:exhaustion} supplies every irreducible summand.
It is therefore essentially surjective onto
\(\operatorname{Rep}^{f}(A_V)\).  Full faithfulness, the unitary tensor
structure and compatibility with the DHR braiding were established in
Theorem~\ref{thm:cwx-embedding}, so the functor is the asserted
equivalence.

For \(V=\mathbb C\Omega\), the scalar net has \(\mu\)-index one and
its additive dualizable representation category is the category of
finite-dimensional Hilbert multiplicity spaces.  The same is true of
finite-length unitary \(V\)-modules, with their ordinary tensor product
and flip braiding.  This gives the stated conclusion in the remaining
case as well.
\end{proof}

Together with the energy bounds and strong locality of
Theorem~\ref{thm:energy} and Corollary~\ref{cor:locality}, the
unitarization and positive fusion forms of
Theorems~\ref{thm:strong-unitarity} and~\ref{thm:positive-fusion}, and
the all-module integration of Theorem~\ref{thm:integration}, this proves
all five assertions of Theorem~\ref{thm:main}.  Complete unitarity
identifies \(\operatorname{Rep}(V)\) with
\(\operatorname{Rep}^{u}(V)\); the equivalence just obtained is the
specified CWX functor on this category, including finite direct sums.

The equivalence also gives a formula for the two-interval index.
Let $M_1,\ldots,M_s$ represent the simple $V$-modules, with positive
forms, and write $d(M_i)$ for their positive intrinsic categorical
dimensions in $\Rep^u(V)$. A unitary tensor equivalence preserves these
dimensions, so they equal the statistical dimensions of the corresponding
net sectors. The two-interval index formula
\cite[Theorem~33]{KLM} therefore gives
\[
 \mu(A_V)=\sum_{i=1}^s d(M_i)^2.
\]
Thus the index of the net is the global dimension of the VOA's unitary
fusion category.

\subsection{Finite-index local extensions}
\label{sec:extensions}

Gui's comparison theorems now give an extension-level application of the
base equivalence. We use the separable-representation and normalized
extension conventions of
\cite[Definitions~3.9 and~5.1]{GuiExtensions}: the VOA embedding is
isometric and preserves the vacuum and conformal vector, while the net
extension includes the specified isometric vacuum embedding and
compatibility of the projective diffeomorphism actions.

\begin{corollary}[Finite-index extension correspondence]
\label{cor:finite-index-extensions}
Let $V$ be a simple unitary strongly rational VOA. Up to isomorphism
over the embedded base, $U\mapsto A_U$ gives a one-to-one correspondence
between the following classes:
\begin{enumerate}[label=\textup{(\arabic*)}]
\item normalized CFT-type unitary VOA extensions $V\hookrightarrow U$
whose underlying $V$-module is grading-restricted and of finite length;
\item normalized irreducible local conformal-net extensions of $A_V$
of finite index.
\end{enumerate}
The VOA class may equivalently be described by simple CFT-type algebraic
extensions with the same conformal vector and finite-length underlying
$V$-module, equipped with their unique unitary structure extending that
of $V$.

Let $P$ be the normalized haploid commutative $C^*$-Frobenius algebra
of $U$ in $\Rep^u(V)$, and let $\Theta$ be its transport to
$\Rep^f(A_V)$ under the base equivalence. The Q-system extension
$B_\Theta$ equals $A_U$ as a conformal net on $\mathcal H_U$.
Moreover, $A_U$ is completely rational and
\[
 F^U_{\mathrm{CWX}}\circ F_{\mathrm{VOA}}
   =F_{\mathrm{CN}}\circ\widetilde F^V_{\mathrm{CWX}}
\]
as braided unitary tensor functors with Gui's tensorators. Here
$\widetilde F^V_{\mathrm{CWX}}$ transports unitary local (dyslectic)
$P$-modules to local $\Theta$-modules, and $F_{\mathrm{VOA}}$ and
$F_{\mathrm{CN}}$ are the canonical equivalences from these categories
to $\Rep^u(U)$ and $\Rep^f(A_U)$, respectively. In particular,
$F^U_{\mathrm{CWX}}$ is a braided unitary tensor equivalence onto all
finite-index sectors of $A_U$.
\end{corollary}

\begin{proof}
If $V=\mathbb C\Omega$, the common zero conformal vector and CFT type
force $U=\mathbb C\Omega$. A finite-index extension of the scalar net
has finite-dimensional interval algebras. Since the interval factors of
a nontrivial irreducible conformal net are type~III
\cite[Section~3.3, property~(7)]{CKLW}, this extension is also scalar.
The functors and tensorators then reduce to those of finite-dimensional
Hilbert spaces. We may therefore assume $V\ne\mathbb C\Omega$.

Theorems~\ref{thm:positive-fusion} and~\ref{thm:strong-unitarity}
together give complete unitarity, and
Corollary~\ref{cor:charged-general} gives energy bounds and strong
intertwining for every irreducible unitary intertwining operator.
These are Gui's Condition~I
\cite[Definition~4.15]{GuiExtensions}. Theorem~\ref{thm:rationality}
adds complete rationality and essential surjectivity of the CWX functor,
so $V$ satisfies Condition~I+ of
\cite[Definition~6.17]{GuiExtensions}. We use Gui's Wassermann tensorator
on the CWX functor; its underlying functor is the one identified in
Section~\ref{sec:comparison}.

Gui's Theorems~3.13 and~5.3 classify the stated normalized net and unitary
VOA extensions by normalized haploid commutative $C^*$-Frobenius algebras
\cite{GuiExtensions}. A finite-index net extension has a dualizable
Q-system object, hence one in $\Rep^f(A_V)$; every such object is a finite
sum of VOA completions by Theorem~\ref{thm:rationality}, so is separable
as required by Gui's convention. The braided unitary base equivalence
therefore transports all these algebra objects and their isomorphisms.
For the algebraic VOA formulation, regularity from
Section~\ref{sec:preliminaries} and complete unitarity allow
\cite[Theorem~4.7]{CGGH}, which gives the asserted unique unitary
structure for simple CFT-type extensions.

Finally, Gui's Corollary~6.5 and Theorem~6.18(2) pass Condition~I+ to
$U$. His Theorem~6.11 and Corollary~6.16 then give $A_U=B_\Theta$ and
the displayed equality of braided unitary tensor functors, with the full
finite-index sector category as target \cite{GuiExtensions}.
\end{proof}
